\subsection{Discussion}
\todo[inline]{Discuss the results and the whole process and stuff like that}

\todo[inline]{Also shoudl explain that the local map generation algorithm is working well in real time and has improvements to it all, that it can run even on a RPI4 and that there was a lot of effort to make the algorithm scalable and easily configurable on a ny drone with just plug and  play library an dexposed tuning varaibles and documentation as well as ROS2 package to try and deply directly on any robot with sonar. NOTE that the algoruthm is still bruite force and fine tuning is needed to run real time to not owerwhelme mainly the pruning parameters and map resolution. Never the less is good and runs well comparable to other modern day methods of sonar mapping give examples. However say that in order for it to go beyond the modern methods and exeed it it has potential, mainly i paralelliability as each sample is independent and reinds a lot of how computer graphics do ilumination of diffrenet materials, could make use of computer graohics concepts and GPU parallel processes to make shader code and convert teh code to GPU interactions to scale teh algorithm even more, thsi soudl mean that even the an SBC with a integrated graphics card could run massives maps with ultra high end resoliution sonars in an acuatte way. SO since we have shwos that thsi algorithms works well on embedded hardware and runs real time and gives good maps and results, it can be asumed the same woudl be true for the GPU and that it woudl enable generating massive maps with low overhead and easy integration with existing systems jesjes}

\todo[inline]{discuss dead reconing in map and that loop closure is important. EXPLAIN THAT THE MAP IS STABLE LOCALLY AND using it locally is the main purpose to extra t stable features that then can be used for DA and loop closure to correct state estimate so it doesn't dead recon and make even more stable maps, a pisitive feedback -> then when we build global map we can use these local map mosaics to reconstruct a much more stable result global full map then without slam as seeb in result map generation where loop closure was not included so dead reconing effect over long distances, link rhe picture from results jesjes map generation}

\todo[inline]{Talk about feature extraction and conpression of data. basically this: duscuss that the main point is to extract features and their probabilities/characteristics so that instead of using map that overime changes and is super big and slow to mathch features on a map, you create a database with ALL global features, this way when doing data association you don't need to use the whole map that is expensive you just check the weights and characteristics. if they match well then it is probable they are candidates for loop closures. this way you compress data of the map and build a database/global map of features and their probabalistic weights, and that is MUCH faster, i stead of like 2GB data to process each time you want to associate features and use the map, you use the feature probabilistic weights that are like 1MB in size and can be checked really fast and globally and pruned for loop closure candidates}

\todo[inline]{also extra, since we are using bounded boxes we could also merge multiple objects tgat are close or overlapping compressing the landmark data even more for more performance jesjes}

\todo[inline]{another thing. even though its real time, there are still improvements available. A lot of the run time is used to move large amounts of data, and as we have seen in results from map generation performance picture BOTH figures ref, lower memory overhead drastically increases runntime. So a proposal would be to use LOSSLESS compression algorithms to drastically decrease memory overhead while preserving information. then one can work on compressed map instead as it will alter information in same way as if we did original uncompressed, for example inage processing would also become faster, but bigest benefit is that we would not need to pass as much data, witch was the main runtime consumer, sending large data amounts jesjes.}

\todo[inline]{discuss map generation can be optimized even further as most of this stuuf is computer graphics and earily similar to rasterization and ilumination techniquise used, so mass pararell GPU acceleration would make this even nore edicient, enabling small compute budgets ro maintain not only small to medium sized map, butteully massive maps that spans GB of data without issue, making it even more efficient and possible to use high fidelity sensors where you only need a single core on a single SBC to maintain the sytem instead of 2-4 depending on map fidelity, while other cores can be used for DA, SLAM and guidance systems.
NOTE: today's new implementation can already do real time while operating other tasks, but it would be even more efficient with GPU acceleration.}

\todo[inline]{EXPLAIN the performance of fearure extreaction explain that the height estimation is the worst time consumer with 40percent of all runtime, Filtering at the start is 20percent of it all and semantic segmentation is also around 20percent of it all this means the rest is just 20percent of the consumption these values are very rought estimates after profiling. So jeeeeee there is space for improvement however it is real time jesjes ++++ ALSO DISCUSS at the same time that height estimation could be done a lot better right now it is a very rough estimate.}



